\FloatBarrier
\section{结论}

本文研究跨镜头切换的流式重建问题，其中递归历史具有相互冲突的作用：它为建立相邻镜头之间的关系提供必要参照，但持续传播同一状态又可能在新镜头内引起随时间变化的漂移。我们的核心观察是，历史应被临时读取以完成镜头间对齐，而不应作为新镜头的递归状态继续传播。基于这一观察，\method{} 通过切换前历史推断镜头间关系，同时仅允许独立重新初始化的状态驱动后续重建。它将各镜头中的相机、人体和场景几何表示到统一世界坐标系，并维持跨镜头人物身份，同时按照输入顺序处理每一帧，无需访问未来帧或进行逐视频优化。三个多人体基准上的实验表明，\method{} 能够持续改善统一坐标系下的人体重建、相机轨迹估计和人物身份关联，并在大幅视角变化下取得最明显的收益。如何在包含复杂剪辑、人物反复进出以及快速视觉变化的长视频中保持稳定重建，仍是值得进一步研究的问题。
